\section*{Bab \ref{ch:g8:speed} --- Perbandingan Senilai, Kecepatan dan Rata-rata}

\begin{solution}{exo:g8:speed:1}
$\frac{150}{2} = 75$ km/jam. \quad
$45$ menit $= 0.75$ jam: $\frac{27}{0.75} = 36$ km/jam. \quad
$\frac{100}{12.5} = 8$ m/detik.
\end{solution}

\begin{solution}{exo:g8:speed:2}
$2 \times 85 = 170$ km. \quad
$40$ menit $= \frac23$ jam: $90 \times \frac23 = 60$ km. \quad
$\frac{315}{90} = 3.5$ jam $= 3$ jam $30$ menit.
\end{solution}

\begin{solution}{exo:g8:speed:3}
$72 \div 3.6 = 20$ m/detik. \quad
$5 \times 3.6 = 18$ km/jam. \quad
$108 \div 3.6 = 30$ m/detik.
\end{solution}

\begin{solution}{exo:g8:speed:4}
$600 \div 15 = 40$ menit. \quad
$2$ jam $30 = 150$ menit: $15 \times 150 = 2250$ L.
\end{solution}

\begin{solution}{exo:g8:speed:5}
$\frac{3.30}{1.2} = 2.75$ euro/kg melawan
$\frac{2.16}{0.8} = 2.70$ euro/kg: tawaran yang kedua (sedikit) lebih
menguntungkan.
\end{solution}

\begin{solution}{exo:g8:speed:6}
\[
\bar v = \frac{25 \times 12.4 + 35 \times 10}{60}
= \frac{310 + 350}{60} = \frac{660}{60} = 11 .
\]
\end{solution}

\begin{solution}{exo:g8:speed:7}
$d = 340 \times 6 = 2040$ m $\approx 2$ km.
\end{solution}

\begin{solution}{exo:g8:speed:8}
Jaraknya: $2 \times 5 = 10$ km, lalu $1.5 \times 4 = 6$ km:
seluruhnya $16$ km. Waktu yang berlalu: $2 + 0.5 + 1.5 = 4$ jam.
\omterm{def:g8:speed:def}{Kecepatan rata-ratanya}: $\frac{16}{4} = 4$ km/jam.
\end{solution}

\begin{solution}{exo:g8:speed:9}
Berbobot:
$\frac{3 \times 15 + 2 \times 9 + 1 \times 12}{3 + 2 + 1}
= \frac{45 + 18 + 12}{6} = \frac{75}{6} = 12.5$.
Tanpa bobot: $\frac{15 + 9 + 12}{3} = 12$. Keduanya berbeda karena
koefisiennya memberi nilai $15$ bobot tiga kali bobot nilai
$12$.
\end{solution}

\begin{solution}{exo:g8:speed:10}
\emph{1.} Waktunya: $\frac{120}{80} = 1.5$ jam dan
$\frac{120}{120} = 1$ jam. \omterm{def:g8:speed:def}{Kecepatan rata-ratanya}:
$\frac{240}{2.5} = 96$ km/jam.

\emph{2.} Mobilnya menghabiskan \emph{lebih banyak waktu} pada
$80$ km/jam ($1.5$ jam) daripada pada $120$ km/jam ($1$ jam): jadi
\omterm{def:g8:speed:def}{kecepatan} yang lambat berbobot lebih besar pada \omterm{def:g8:speed:average}{rata-rata berbobot}
waktunya, sehingga menariknya ke bawah titik tengah $100$.
\end{solution}

\begin{solution}{exo:g8:speed:11}
Celah $18$ km menutup pada $5 + 4 = 9$ km/jam: mereka bertemu sesudah
$\frac{18}{9} = 2$ jam, yaitu pukul $12{:}00$. Ana ketika itu sudah
berjalan $2 \times 5 = 10$ km: mereka bertemu $10$ km dari desa Ana
(dan $8$ km dari desa Bayu --- periksa: $10 + 8 = 18$).
\end{solution}

\begin{solution}{pb:g8:speed:1}
\textbf{1.} Menurut \cref{def:g8:speed:def}, $t_1 = \frac{d}{v_1}$ dan
$t_2 = \frac{d}{v_2}$; seluruh perjalanannya berlangsung
$\frac{d}{v_1} + \frac{d}{v_2}$ untuk jarak $2d$.

\textbf{2.} Faktorkan $d$ keluar dari \omterm{def:g4:fractions:def}{penyebutnya} lalu sederhanakan
sampai lenyap:
\[
\bar v
= \frac{2d}{d\left(\frac{1}{v_1} + \frac{1}{v_2}\right)}
= \frac{2}{\ \frac{1}{v_1} + \frac{1}{v_2}\ } .
\]
\omterm{def:g4:fractions:def}{Penyebut} bersama di dalamnya:
$\frac{1}{v_1} + \frac{1}{v_2} = \frac{v_2 + v_1}{v_1 v_2}$, dan
membagi berarti mengalikan dengan \omterm{def:g8:fractions:inverse}{kebalikannya}
(\cref{thm:g8:fractions:div}):
\[
\bar v = 2 \times \frac{v_1 v_2}{v_1 + v_2}
= \frac{2\, v_1 v_2}{v_1 + v_2} .
\]

\textbf{3.} $\dfrac{2 \times 30 \times 15}{30 + 15}
= \dfrac{900}{45} = 20$ km/jam, seperti pada
\cref{ex:g8:speed:twolegs}; dan
$\dfrac{2 \times 80 \times 120}{80 + 120} = \dfrac{19\,200}{200}
= 96$ km/jam, seperti yang ditemukan pada \cref{exo:g8:speed:10}.

\textbf{4.} Rata-ratanya hanya bergantung pada kedua \omterm{def:g8:speed:def}{kecepatannya},
bukan pada $d$: perjalanan pulang pergi $1$ km dan perjalanan pulang
pergi $100$ km pada kedua \omterm{def:g8:speed:def}{kecepatan} yang sama punya \omterm{def:g8:speed:def}{kecepatan}
rata-rata yang persis sama. (Itulah sebabnya rumusnya dapat diperiksa
pada perjalanan yang panjangnya berbeda-beda.)

\textbf{5.} Dalam waktu $T$ pada tiap \omterm{def:g8:speed:def}{kecepatannya}, jaraknya adalah
$v_1 T$ dan $v_2 T$, jadi
\[
\bar v = \frac{v_1 T + v_2 T}{2T}
= \frac{(v_1 + v_2)\,T}{2T}
= \frac{v_1 + v_2}{2} .
\]
Rata-rata \omterm{def:g8:speed:def}{kecepatan} selalu dibobot menurut \emph{waktu}
(\cref{def:g8:speed:average}). Dengan waktu yang sama kedua
\omterm{def:g8:speed:def}{kecepatannya} membawa bobot yang sama: jadi titik tengah biasa. Dengan
jarak yang sama, \omterm{def:g8:speed:def}{kecepatan} yang lebih lambat menempati waktu yang
\emph{lebih} banyak, jadi bobotnya lebih besar --- dan rata-ratanya
menggeser ke arahnya.

\textbf{6.} Untuk $(30, 15)$: $H = 20$ dan $A = 22.5$. Untuk
$(80, 120)$: $H = 96$ dan $A = 100$. Rata-rata aritmetikanya menang
kedua kalinya.

\textbf{7.} Pada \omterm{def:g4:fractions:def}{penyebut} bersama $2(v_1 + v_2)$:
\[
A - H
= \frac{(v_1 + v_2)^2 - 4\, v_1 v_2}{2\,(v_1 + v_2)} ,
\]
dan \omterm{def:g4:fractions:def}{pembilangnya} menjabar dengan kesamaan pada
\cref{pb:g8:equations:1}:
\[
(v_1 + v_2)^2 - 4 v_1 v_2
= v_1^2 + 2 v_1 v_2 + v_2^2 - 4 v_1 v_2
= v_1^2 - 2 v_1 v_2 + v_2^2
= (v_1 - v_2)^2 .
\]

\textbf{8.} Kuadrat tidak pernah negatif, dan $2(v_1 + v_2)$ bernilai
positif: jadi $A - H \geq 0$, yaitu $H \leq A$, untuk setiap pasangan
\omterm{def:g8:speed:def}{kecepatan} yang positif. Kesamaannya menuntut $(v_1 - v_2)^2 = 0$,
yaitu $v_1 = v_2$: begitu kedua \omterm{def:g8:speed:def}{kecepatannya} berbeda, rata-rata
perjalanan pulang perginya turun benar-benar di bawah titik
tengahnya --- kejutannya adalah sebuah teorema.

\textbf{9.} Kalikan:
\[
H \times A
= \frac{2\, v_1 v_2}{v_1 + v_2} \times \frac{v_1 + v_2}{2}
= v_1 v_2 ,
\]
sesudah mencoret faktor $2$ dan faktor $v_1 + v_2$.

\textbf{10.} Dengan rumusnya: $\bar v = \dfrac{2 \times 40 \times 60}
{40 + 60} = \dfrac{4800}{100} = 48$ km/jam. Lewat jalan panjangnya:
$\frac{60}{40} = 1.5$ jam dan $\frac{60}{60} = 1$ jam, jadi $120$ km
dalam $2.5$ jam: $\frac{120}{2.5} = 48$ km/jam. Keduanya sesuai.

\textbf{11.} Rata-rata $30$ km/jam sepanjang $30$ km berarti waktu
seluruhnya paling lama $\frac{30}{30} = 1$ jam. Sedangkan $15$ km
pertamanya pada $15$ km/jam sudah makan waktu $\frac{15}{15} = 1$ jam:
\emph{seluruh} jatah waktunya terpakai, dengan separuh lintasannya
belum ditempuh.

\textbf{12.} Pada $90$ km/jam separuh keduanya makan waktu
$\frac{15}{90} = \frac16$ jam, jadi lombanya makan $\frac76$ jam dan
\[
\bar v = \frac{30}{\ 7/6\ } = \frac{180}{7} \approx 25.7
\text{ km/jam} < 30 .
\]
Berapa pun \omterm{def:g8:speed:def}{kecepatan} separuh keduanya, lamanya adalah bilangan yang
positif, jadi waktu seluruhnya melampaui $1$ jam dan rata-ratanya
tetap benar-benar di bawah $30$ km/jam. (Dengan rumusnya:
$\frac{2 \times 15 \times v_2}{15 + v_2} = 30$ akan memaksa
$30 v_2 = 450 + 30 v_2$, yaitu $0 = 450$ --- tidak ada \omterm{def:g8:speed:def}{kecepatan}
$v_2$ yang berhasil.)

\textbf{13.} Rata-rata $20$ km/jam mengizinkan $\frac{30}{20} = 1.5$
jam; sesudah jam pertamanya, $15$ km tersisa untuk $0.5$ jam: jadi ia
perlu $\frac{15}{0.5} = 30$ km/jam. Periksa:
$H(15, 30) = \frac{2 \times 15 \times 30}{45} = \frac{900}{45}
= 20$ km/jam.

\textbf{14.} Tiap menitnya, kerannya mengisi $\frac{1}{20}$ dan
$\frac{1}{30}$ bagian baknya; bersama-sama
\[
\frac{1}{20} + \frac{1}{30}
= \frac{3}{60} + \frac{2}{60}
= \frac{5}{60} = \frac{1}{12}
\]
bagian baknya tiap menit: jadi baknya terisi dalam $12$ menit. Dan
$H(20, 30) = \frac{2 \times 600}{50} = 24 = 2 \times 12$: kedua
kerannya bersama-sama makan waktu persis \emph{setengah \omterm{pb:g8:speed:1}{rata-rata
harmonik}} waktu masing-masingnya --- laju yang berjumlah, jadi ini
wilayah \omterm{pb:g8:speed:1}{rata-rata harmonik} lagi.

\textbf{15.} Jaraknya sama pada tiap arahnya, jadi
\[
\bar v = \frac{2 \times 900 \times 700}{900 + 700}
= \frac{1\,260\,000}{1\,600} = 787.5 \text{ km/jam},
\]
di bawah $800$ km/jam pada udara yang tenang. Dalam satu kalimat:
pesawatnya menghabiskan waktu \emph{lebih lama} terbang melawan angin
daripada bersamanya, jadi ruas yang lambat memperoleh bobot yang
lebih besar pada \omterm{def:g8:speed:average}{rata-rata berbobot} waktunya --- angin hanya dapat
memperpanjang perjalanan pulang pergi.

\end{solution}
